% ============================================================
% Come si usano, argomento 12: PDF condizionate e funzioni di variabili continue
% ============================================================
\argomento{PDF condizionate e funzioni di variabili continue}{Formule a p.~\pageref*{f:pdfcond}}[PDF condizionate
  e funzioni]

\begin{scelta}{Come scegliere la formula}
\renewcommand{\arraystretch}{1.25}
\begin{tabular}{@{}P{0.55\linewidth}P{0.39\linewidth}@{}}
  \toprule
  \textbf{Se il testo\ldots} & \textbf{si usa\ldots} \\
  \midrule
  dice che $X$ è caduta in un intervallo, o dà un altro evento & la densità condizionata \vv{pdfcond} \\
  fa dipendere $X$ da un caso che non si conosce (quale fornitore, quale canale) & la probabilità totale
    \vv{totalepdf} \\
  chiede la densità di $Y = g(X)$, con $g$ monotona & \vv{monotona} \\
  chiede la densità di $Y = g(X)$, con $g$ che torna indietro ($x^2$, coseno) & la somma sulle soluzioni
    \vv{nonmonotona} \\
  chiede solo la media di $Y = g(X)$ & il teorema fondamentale (argomento 13): la densità di $Y$ non serve \\
  chiede come generare numeri con una densità data & il metodo dell'inversione \vv{inversione} \\
  quantizza $X$ su un numero finito di livelli & la conversione A/D \vv{conversionead} \\
  \bottomrule
\end{tabular}

Schema per $f_Y$: (1) si trova l'alfabeto di $Y$; (2) per ogni $y$ si risolve $g(x) = y$; (3) si calcola
$\abs{\dot g}$ nelle soluzioni; (4) si sommano i rapporti $f_X / \abs{\dot g}$, scritti in funzione di $y$.
\end{scelta}

\scheda{pdfcond}{
  \usoquando{Si sa che $X$ è caduta in un intervallo (o che si è verificato un altro evento), e serve la sua densità
    con questa informazione.}
  \usocome{Si taglia la densità fuori dall'intervallo e la si «rialza» dividendola per la probabilità dell'intervallo,
    così che l'area torni a 1.}
  \usoesempio{$X \sim \mathcal{E}(1)$, sapendo che $X \leq 1$: $f_{X \mid A}(x) = \frac{e^{-x}}{1 - e^{-1}} \simeq
    \frac{e^{-x}}{0{,}632}$ su $[0, 1]$. Una Laplace con $\lambda = 1$ ristretta a $[-1, 2]$ va divisa per
    $F_X(2) - F_X(-1) \simeq 0{,}748$.}
}

\scheda{totalepdf}{
  \usoquando{$X$ dipende da un caso che non si conosce, e noto il caso la sua densità è nota: la densità è una
    «miscela».}
  \usocome{Si pesano le densità (o le CDF, o le medie) dei casi con le loro probabilità.}
  \usoesempio{Componenti di due fornitori in parti uguali, durate esponenziali di media 1 e 3 anni:
    $f_X(x) = \frac{1}{2}\, e^{-x} u(x) + \frac{1}{2} \cdot \frac{1}{3}\, e^{-x/3} u(x)$ ed
    $E[X] = \frac{1}{2} \cdot 1 + \frac{1}{2} \cdot 3 = 2$ anni, senza integrali.}
}

\scheda{monotona}{
  \usoquando{$Y = g(X)$ con $g$ sempre crescente o sempre decrescente sul supporto di $X$: un solo $x$ per ogni $y$.}
  \usocome{Si ricava $x(y)$ da $y = g(x)$, si calcola $\dot g$ in $x(y)$ e si divide. Dove $g$ è ripida la densità si
    abbassa, dove è piatta si alza: la probabilità si conserva.}
  \usoattenzione{Al denominatore va il \emph{valore assoluto} della derivata, e tutto va scritto in funzione di $y$.}
  \usoesempio{$X \sim \mathcal{U}(0, 1)$, $Y = 3X + 2$: $f_Y(y) = \frac{1}{3}$ su $(2, 5)$, cioè $Y \sim
    \mathcal{U}(2, 5)$. Il faro: $\Theta \sim \mathcal{U}\left(-\frac{\pi}{2}, \frac{\pi}{2}\right)$ e
    $X = \tan\Theta$; con $\dot g = 1 + x^2$ si trova $f_X(x) = \frac{1}{\pi}\, \frac{1}{1 + x^2}$, una Cauchy.}
}

\scheda{nonmonotona}{
  \usoquando{$Y = g(X)$ con $g$ che torna indietro, come $x^2$ o una sinusoide: lo stesso $y$ viene da più $x$.}
  \usocome{Si trovano tutte le soluzioni di $g(x) = y$ e si somma un termine $f_X / \abs{\dot g}$ per ciascuna.}
  \usoattenzione{Una densità ottenuta così può tendere a infinito (dove $\dot g \to 0$) e restare una densità: conta
    l'area, non l'altezza.}
  \usoesempio{$X \sim \mathcal{U}(-1, 1)$ e $Y = X^2$: due soluzioni $\pm\sqrt{y}$, ciascuna con densità
    $\frac{1}{2}$, quindi $f_Y(y) = \frac{1/2 + 1/2}{2\sqrt{y}} = \frac{1}{2\sqrt{y}}$ su $(0, 1)$. Con
    $X \sim \mathcal{C}(0, 1)$ si trova $f_Y(y) = \frac{1}{\pi \sqrt{y}\,(1 + y)}\, u(y)$.}
}

\scheda{cdfuniforme}{
  \usoquando{Per capire il metodo dell'inversione: qualunque variabile continua, passata nella propria CDF, diventa
    uniforme fra 0 e 1.}
  \usoesempio{$X \sim \mathcal{E}(\lambda)$: $Y = 1 - e^{-\lambda X}$ è uniforme su $(0, 1)$.}
}

\scheda{inversione}{
  \usoquando{Si deve simulare una variabile con una densità data, partendo da numeri uniformi in $(0, 1)$ (quelli che
    produce un computer).}
  \usocome{Si risolve $u = F_X(x)$ rispetto a $x$: la formula trovata, applicata a $U$, dà la variabile voluta.}
  \usoesempio{Esponenziale: da $u = 1 - e^{-\lambda x}$ si ricava $x = -\frac{1}{\lambda}\ln(1 - u)$. Con
    $u = 0{,}5$ si ottiene $z = \frac{\ln 2}{\lambda} \simeq \frac{0{,}693}{\lambda}$.}
}

\scheda{arcoseno}{
  \usoquando{Si osserva una sinusoide in un istante a caso (o con una fase a caso): il valore letto ha questa
    densità.}
  \usocome{È un caso di \vv{nonmonotona}: due soluzioni per periodo, e $\abs{\sin} = \sqrt{1 - \cos^2}$ per scrivere
    $\abs{\dot g}$ in funzione di $y$.}
  \usoesempio{Con $A = 1$ la densità vale $\frac{1}{\pi} \simeq 0{,}318$ in $y = 0$ e tende a infinito verso $\pm 1$: la
    sinusoide passa più tempo vicino ai massimi e ai minimi, dove rallenta.}
}

\scheda{conversionead}{
  \usoquando{Una variabile continua viene rappresentata con un numero finito di livelli (un convertitore A/D, una
    quantizzazione a $R$ bit).}
  \usocome{La probabilità di ogni livello è la probabilità del suo intervallo: una differenza di CDF.}
  \usoesempio{$X \sim \mathcal{L}(1)$, 3 bit, soglie $-\infty, -3, \ldots, 3, +\infty$: $p_Y(0{,}5) = F_X(1) - F_X(0) =
    \frac{1}{2} - \frac{1}{2}\, e^{-1} \simeq 0{,}316$ e $p_Y(3{,}5) = \frac{1}{2}\, e^{-3} \simeq 0{,}025$.}
}
